\section{Causal interpretation and benchmark equivalence}

The population target in \cref{obj:def:model,obj:def:null} uses continuous representatives of conditional means. Its interpretation therefore requires both invariance to conditional-distribution versions and a causal representation that reproduces the observed law. The following result establishes these properties and supplies legal constant-effect and heterogeneous-effect witnesses for every admissible parameter tuple.

\begin{lemmav}{lem:causal-nonempty}[Nonemptiness and causal identification]
Fix an admissible parameter tuple \(v=(p,\alpha,\beta,\gamma)\) as in \cref{obj:def:model}. Let \(\lambda\) denote Lebesgue probability measure on \([0,1]\), and define the heterogeneity distance, its model supremum, and the witness distance by
\[
d(P)=
\left\{\int_0^1
\left(\tau_P(x)-\int_0^1\tau_P(z)\,d\lambda(z)\right)^2
\,d\lambda(x)\right\}^{1/2},
\qquad
D_v=\sup_{P\in\mathcal M_v}d(P),
\qquad
d_0=\frac{1}{8\sqrt{2}}.
\]
The set \(\{d(P):P\in\mathcal M_v\}\) is bounded above, and
\[
d_0\le D_v\le\frac12.
\]

For every \(P\in\mathcal M_v\), with model membership understood as in \cref{obj:def:model}, the continuous propensity, baseline mean, and contrast representatives are determined uniquely by the observed law: every other observational representation of the same probability law has exactly the same \(e_P\), \(m_{0,P}\), and \(\tau_P\), pointwise on \([0,1]\). Consequently, the treated mean representative \(m_{1,P}=m_{0,P}+\tau_P\) is also uniquely determined, and
\[
|m_{1,P}(x)-m_{1,P}(z)|
\le 40|x-z|^{\min\{\beta,\gamma\}}
\qquad (x,z\in[0,1]).
\]

In the completion \(\mathbb Q_P\) of \cref{obj:def:causal-completion}, let \(Y(0)\) and \(Y(1)\) denote the baseline and treated potential-outcome coordinates. The distribution under \(\mathbb Q_P\) of
\[
\bigl(X,A,(1-A)Y(0)+AY(1)\bigr)
\]
is \(P\), and
\[
(Y(0),Y(1))\perp A\mid X.
\]
Both potential outcomes are integrable under \(\mathbb Q_P\), and
\[
\mathbb E_{\mathbb Q_P}\!\left[Y(1)-Y(0)\mid X\right]
=\tau_P(X)
\qquad \mathbb Q_P\text{-almost surely}.
\]
Moreover,
\[
d(P)=0\quad\Longleftrightarrow\quad P\in H_0(v),
\]
where \(H_0(v)\) is defined in \cref{obj:def:null}.

The law with uniform covariate distribution, propensity and mean functions
\[
e_P(x)=\frac12+\frac{\sin(2\pi x)}{10},
\qquad
m_{0,P}(x)=\frac{\cos(2\pi x)}8,
\qquad
\tau_P(x)=\frac{\cos(2\pi x)}8,
\]
and deterministic conditional outcome
\[
Y=m_{0,P}(X)+A\tau_P(X)
\]
belongs to \(\mathcal M_v\) and has \(d(P)=d_0\). Replacing its contrast function by \(\tau_P(x)=1/8\) for every \(x\in[0,1]\), while retaining the displayed propensity, baseline mean, and deterministic conditional-outcome construction, yields a law in \(H_0(v)\).
\end{lemmav}

The uniqueness assertion makes every pointwise restriction on the propensity and mean functions a property of the observed law. The uniform design and overlap conditions in \cref{obj:ass:uniform,obj:ass:overlap} give the conditional arm distributions a common Lebesgue almost-everywhere interpretation; continuity makes their canonical mean representatives unambiguous across the whole interval. In particular, the treated mean inherits the regularity stated in \cref{obj:lem:causal-nonempty} from the separate baseline and contrast restrictions.

Version invariance also applies to the potential-outcome completion. As specified in \cref{obj:def:causal-completion}, changing either conditional outcome distribution on a Lebesgue-null set leaves \(\mathbb Q_P\) unchanged. The completion combines the observed arm distributions with conditionally independent assignment and reproduces the original-record marginal. Its integrable potential-outcome difference has conditional mean \(\tau_P(X)\), so \(d(P)\) measures variation in conditional average treatment effects under this completion. The equivalence between zero distance and membership in \cref{obj:def:null} connects the population criterion to the pointwise constant-effect restriction.

The two witnesses in \cref{obj:lem:causal-nonempty} give the testing problem a nonempty null and a positive separation domain for every admissible tuple. The heterogeneous witness attains the same fixed distance \(d_0\) throughout the parameter domain, while the constant-contrast witness permits a nonzero common effect. Together with the upper bound on \(D_v\), these properties give the cap in \cref{obj:def:critical-radius} a finite, positive population scale.

\subsection{Mean preservation and the benchmark experiment}

The bounded and signed-binary models in \cref{obj:def:bounded-model,obj:def:binary-model} retain the primitive restrictions on assignment and conditional means. Their comparison in \cref{obj:prop:bounded-submodel-comparison} uses two related operations: a model-level conversion specified by conditional arm means, and a record-level randomization available for bounded observations. At the model level, the converted outcome has probabilities
\[
\Pr(+1\mid X=x,A=a)=\frac{1+m_{a,P}(x)}2,
\qquad
\Pr(-1\mid X=x,A=a)=\frac{1-m_{a,P}(x)}2,
\qquad a\in\{0,1\}.
\]
These probabilities describe the signed-binary conversion in \cref{obj:prop:bounded-submodel-comparison}. It retains the propensity, both conditional arm means, and hence the contrast and its distance from constancy. The same result establishes that the bounded, signed-binary, and original models have equal distance suprema.

For a bounded original record, the conversion can instead use its observed outcome directly. To express this randomization on the full real-outcome test carrier, define the measurable sign map
\[
s(y,t)
=
2\mathbf 1\!\left\{
t\le \frac{1+\max\{-1,\min\{y,1\}\}}2
\right\}-1,
\qquad (y,t)\in\mathbb R\times[0,1].
\]
For \(|y|\le1\), a uniform \(t\) gives the probabilities \((1+y)/2\) and \((1-y)/2\) for the two signs, with conditional mean \(y\). The truncation in this definition specifies the map on the entire original-record carrier. Its restriction to the bounded model is precisely the record-level conversion in \cref{obj:prop:bounded-submodel-comparison}.

Write \(\boldsymbol V=(V_1,\ldots,V_n)\) for independent uniform variables on \([0,1]\), independent of the sample and of the public seed \(U\) in \cref{obj:def:original-record-experiment}. The randomized signed-binary sample is
\[
\mathcal D_n^{\,s}(\boldsymbol V)
=
\bigl((X_i,A_i,s(Y_i,V_i))\bigr)_{i=1}^n.
\]
For a test \(\phi\) on the carrier of \cref{obj:def:testing-risk}, its integrated rejection-probability map is
\[
\widetilde\phi(\mathcal D_n,u)
=
\int_{[0,1]^n}
\phi\!\left(
\bigl((X_i,A_i,s(Y_i,t_i))\bigr)_{i=1}^n,u
\right)
\,\lambda^{\otimes n}(d\boldsymbol t),
\qquad
\boldsymbol t=(t_1,\ldots,t_n).
\]
Thus \(\widetilde\phi\) is a jointly Borel map from
\[
\bigl([0,1]\times\{0,1\}\times\mathbb R\bigr)^n\times[0,1]
\quad\text{to}\quad[0,1].
\]
It uses the original-record experiment and its public seed; the additional sign randomization is averaged into the rejection probability.

For \(P\in\mathcal M_w^{\mathrm b}\), let \(P^{\,s}\) denote its signed-binary conversion from \cref{obj:prop:bounded-submodel-comparison}. The rejection-probability preservation assertion there takes the explicit form
\[
\mathsf R_n(P,\widetilde\phi)
=
\mathsf R_n(P^{\,s},\phi).
\]
This identity expresses the statistical role of mean-preserving randomization: rejection probabilities under each bounded law transfer exactly to those under its signed-binary conversion, including laws in the constant-effect null and laws at a specified heterogeneity distance.

\Cref{obj:prop:bounded-submodel-comparison} consequently gives the exact benchmark identity
\[
B_n^{\mathrm{bin}}(w,r)=B_n^{\mathrm b}(w,r),
\qquad n\ge2,\quad 0<r<D_v,
\]
for the risks in \cref{obj:def:binary-risk,obj:def:bounded-risk}. With the common distance supremum and the capping conventions in \cref{obj:def:binary-radius,obj:def:bounded-radius}, the same comparison gives equal capped critical radii for every \(n\ge2\). The signed-binary benchmark therefore represents the bounded mean-effect testing problem at the level of attainable distances, minimax risks, and critical radii.
